%% methodology.tex
%%

%% ==============================
\section{Methodology}
\label{sec:Methodology}
%% ==============================

\subsection{Synthetic Benchmark Generation}
For each of 5 seeds (42-46), 1,200 IaC snippets are generated (\texttt{src/data/iac\_dataset.py}), each independently choosing an easy or hard pattern (50/50) and a safe/misconfigured label (50/50), with 2-4 unrelated filler configuration lines added as noise. Two paired renderings are produced from the identical underlying random choices -- with comments and without -- so that any difference in a model's performance between the two is attributable only to the presence or absence of the comment, not to a different underlying sample.

\subsection{Training and Evaluation}
A 70/30 stratified train/test split is used per seed. The ML detector is trained separately on the rich-context and code-only training sets, then evaluated on the matching test set. The rule-based and hybrid detectors require no training beyond the ML component the hybrid detector wraps (trained on the code-only training set, since that is the setting the improvement targets).

\subsection{Metrics}
Precision, Recall, and F1 -- the same three metrics War et al. use -- computed on the held-out test set for each of: ML Detector (rich context), ML Detector (code-only), Rule-Based Only, and Hybrid Detector (code-only). A confidence-threshold sweep (0.5 to 0.9) for the hybrid detector's ML fallback is also reported, to characterise the precision/recall trade-off rather than a single fixed operating point.

All code and results are in \texttt{iac-security/}.
